% Figure: trapezoidal back-EMF of the three phases of a brushless DC
% machine against electrical angle, with the six-step commutation windows
% marked. The flat top is the design feature: each phase is energised only
% while its back-EMF is at the plateau, so the torque per ampere stays
% nearly constant through the electrical cycle.
\begin{figure}[htbp]
\centering
\begin{tikzpicture}[>=Latex]
\begin{axis}[
  width=12.5cm, height=6.5cm,
  xlabel={electrical angle $\theta_e$ (degrees)},
  ylabel={phase back-EMF / $k_e\omega$},
  xmin=0, xmax=360,
  ymin=-1.35, ymax=1.35,
  xtick={0,60,120,180,240,300,360},
  ytick={-1,-0.5,0,0.5,1},
  legend pos=south east,
  legend cell align=left,
  font=\scriptsize,
  grid=both,
  grid style={enggray!20},
]
% Phase A trapezoid: flat +1 over 30..150, ramp down 150..210, flat -1
% over 210..330, ramp up 330..360 and 0..30.
\addplot[engblue, line width=1.1pt] coordinates {
  (0,1) (30,1) (150,1) (210,-1) (330,-1) (360,1)
};
\addlegendentry{phase A}
% Phase B shifted by +120 deg.
\addplot[engred, line width=1.1pt] coordinates {
  (0,-1) (90,-1) (150,1) (270,1) (330,-1) (360,-1)
};
\addlegendentry{phase B}
% Phase C shifted by +240 deg.
\addplot[engamber, line width=1.1pt] coordinates {
  (0,1) (30,1) (90,-1) (210,-1) (270,1) (360,1)
};
\addlegendentry{phase C}
% shade one commutation window (60 deg sector where A is +1, B is -1)
\addplot[enggray!18, fill=enggray!18, draw=none] coordinates {
  (30,-1.35) (90,-1.35) (90,1.35) (30,1.35)
} \closedcycle;
\node[enggray, anchor=south, font=\scriptsize] at (axis cs:60,-1.32)
  {one $60^{\circ}$ step};
\end{axis}
\end{tikzpicture}
\caption[Trapezoidal back-EMF of a brushless DC machine]{Back-EMF of the
three phases of a brushless DC machine against electrical angle, each
normalised by the back-EMF constant times speed, $k_e\omega$. The
trapezoidal shape is the design target of a concentrated-winding machine:
a flat plateau of width $120^{\circ}$ electrical on each phase, joined by
$60^{\circ}$ ramps. The drive energises the two phases whose back-EMFs sit
on their plateaus during each $60^{\circ}$ commutation step (one window is
shaded), so the product of phase current and back-EMF, and therefore the
torque, stays nearly constant through the electrical cycle. The torque
ripple of the real machine is the departure of the actual waveform from
this idealised trapezoid.}
\label{fig:vol04ch08:bldc-backemf}
\end{figure}
